\section{Related work and a defensible contribution boundary}
This is a scoped primary-source review, not an exhaustive novelty search.
The bibliography identifies the versions inspected. Documentation establishes
described workflows; it does not establish a competing system's error rate,
the absence of an internal technique or Texloom's superiority.

\subsection{Synchronization as evidence}
Laurens' author-maintained SyncTeX API documents source/display queries,
multiple query results, weak matching, line fallback when an exact column is
unavailable and ambiguity involving PDF forms~\cite{synctex}. These are useful
navigation facilities, not a specification for correct source updates.
Texloom therefore uses synchronization as a candidate-generating layer and
requires additional source and output evidence before committing an edit.
The inspected API source is pinned to a commit. The 2008 TUGboat full text
could not be retrieved reliably during this review, so the report does not
claim a full-text review of that article.

\subsection{Graphical scientific and LaTeX editing}
LyX describes a graphical, structure-oriented document processor using
TeX/LaTeX~\cite{lyx}. TeXmacs describes structured scientific editing with
interactive high-quality typesetting~\cite{texmacs}. These systems already
address the broad interface goal of editing scientific documents visually.
Differences in document representation, source ownership and supported edits
must be stated before any comparison.

The archived BaKoMa distribution README describes TeXWord editing LaTeX source
through a rendered-page interface~\cite{bakoma}. It is particularly relevant
prior work against any claim that editing typeset LaTeX pages is itself new.
The inspected text is a historical 2011 distribution description, not evidence
of current availability, internal algorithms or preservation guarantees.
Those questions require additional primary documentation or experiments.

Overleaf documents code/visual editing alongside a compiled PDF preview and
source/PDF navigation~\cite{overleaf}. An interface mode or synchronization
feature is not automatically the displacement contract studied here. Equally,
the documentation's emphasis does not prove that an unmentioned technique is
absent. This report makes workflow distinctions rather than a timing leaderboard.

\subsection{Bidirectional transformations and output-directed programming}
Foster et al. formulate bidirectional transformations using a view and a
putback operation, with round-trip laws~\cite{foster2005}. As a conceptual
analogy, desirable behavior includes leaving a source unchanged when its view
is unchanged and producing an intended new view after a supported update.
Texloom has no proof that arbitrary TeX/PDF behavior satisfies such laws.
Compilation is not generally injective, and the source contains information
that a rendered page cannot recover. Exact Undo is a property of retained
history, not proof of a general bidirectional transformation law. The cited
source is the POPL 2005 conference version, not the expanded 2007 journal paper.

Mayer, Kuncak and Chugh update programs from manipulated output through
bidirectional evaluation and custom lenses~\cite{mayer2018}. Hempel, Lubin and
Chugh combine source and graphical operations in output-directed SVG
programming~\cite{hempel2019}. This is close conceptual prior work: output
manipulation can guide source updates, and provenance constrains the choices.
Texloom instead applies bounded source plans around an existing TeX engine
and checks compiled candidates. It does not implement a reverse evaluator for
general TeX. Any later contribution would need to explain what is learned
about this particular preservation/refusal/cost trade-off beyond assembling
known ideas.

\subsection{Validation and document differences}
Translation validation checks individual translations rather than relying
only on trust in the translating implementation~\cite{pnueli1998}. The
publisher abstract was inspected; the original formal method was not
reproduced here. The useful analogy is checking each candidate. Texloom's
finite source, geometry and raster predicates do not become that paper's
semantic refinement proof.

\emph{latexdiff} marks significant differences between LaTeX
sources~\cite{latexdiff}. Such source review is complementary: it can expose
unrequested edits, but does not by itself establish that unchanged source
rendered identically or that the intended printed occurrence moved. The
evaluation consequently combines declared source changes with independent
rendered evidence and keeps their failure modes separate.

\subsection{Current contribution claim}
The defensible present claim is a documented systems mechanism and a
reproducible way to investigate it in a bounded setting. Potentially useful
findings include which incorrect candidates escape weaker checks, where a
strict contract refuses useful edits, and how much first-selection and
verification cost. No claim of being the first visual LaTeX editor, of formal
verification, of universal preservation or of superior performance follows
from the current audit. A stronger research claim needs a deeper literature
review, reviewed labels and substantive fixed-version comparisons.
